% Papiers de Sophie Germain — Français 9118, lecture modernisée du volume entier.
%
% Pass: Opus 5 (claude-opus-5), 11 September 2026 — first pass, unchecked
% against the views by a human, and an interpretation rather than a record.
% Where this file and the transcriptions differ in substance, the
% transcriptions (batch-01.fr.tex … batch-06.fr.tex) are the record.
%
% Sources read: the six batch transcriptions of this volume and nothing else.
% No image was consulted, and no published edition — not Grun's, not Del
% Centina's, not the printed 1821 « Recherches ». Batches 1–3 were transcribed
% on Fable 5.1, batches 4–6 on Opus 5; batch 6 (views 101–103) is the end of
% the binding and holds no text at all, which is why nothing below draws on it.
%
% What the volume turned out to be. Not a mathematical notebook: the dossier
% of a career as it arrived from outside. Letters received (Cauchy, Delambre,
% Fourier, Legendre, Navier, Poisson, Gauss, Pernety, Villoison, Libri), three
% drafts of her own, a copy of Lagrange's note to the prize commissioners, a
% tutorial in a third hand, a printed offprint, three dated Institut receipts,
% and a dozen loose leaves in her hand. The mathematics is therefore mostly
% other people's, addressed to her — which is exactly what makes it worth
% reading: the leaves record what was said to her, and what she answered.
%
% Two spines carry almost all of it: the vibrating elastic surface (1811–1826)
% and the arithmetic of the Disquisitiones (1804–1808, with a Fermat coda in
% 1819). A third, smaller, runs across both: Euler's vibrating lamina, which
% Legendre and she argue about in four separate pieces.
%
% Notation translated here, each with a footnote at its first use: Lagrange's
% and Legendre's fourth-order operator into the biharmonic; « 1/r + 1/r' »
% into twice the mean curvature and, at small deflection, into the Laplacian;
% Legendre's barred e and vinculum into modern exponents and brackets; Gauss's
% « déterminant » of a form into the determinant of its Gram matrix; his
% (a, b, c) into ax² + 2bxy + cy²; Germain's « résidu » language, where it
% occurs, into congruences.
\documentclass[11pt,a4paper]{article}
% The preamble shared by every transcription of the mathematicians' manuscripts in Gallica.
%
% It rests on one idea: **uncertainty compounds, it does not resolve.** A
% transcription that smooths over a doubtful word has destroyed the only thing
% it was made for — a reader must be able to tell, at every point, what was on
% the page and what was supplied. The apparatus macros below are therefore the
% critical apparatus, not decoration.
%
% The same commands are recognised by `scripts/render.mjs`, which produces the
% site's reading view, and by `scripts/tei.mjs`. A macro added here must be
% added there too, or rendering fails loudly — which is the wanted behaviour.
%
% Comments here are English, like the rest of the repository. The documents
% this preamble sets are French, and that is deliberate: see the skills.

\ProvidesPackage{germain}[2026/09/15 Transcriptions of mathematicians' manuscripts in Gallica]

\usepackage{amsmath,amssymb,amsthm}
\usepackage{mathrsfs}
% Kept from the parent preamble so the renderer's diagram support stays
% available; these manuscripts draw no commutative diagrams, and nothing here uses it.
\usepackage{tikz-cd}
\usepackage[english,french]{babel}
\usepackage{fontspec}

% --- Greek ------------------------------------------------------------------
% The text face stays fontspec's default, Latin Modern, which has no Greek:
% XeTeX drops every glyph it lacks with a line in the log and nothing on the
% page, and the Greek words the manuscripts quote vanished from the PDFs. So
% the Greek and Greek Extended blocks get a XeTeX character class of their own,
% and entering or leaving a run of it switches the family to CMU Serif — the
% same Computer Modern design, with polytonic Greek — and back. Only the
% family changes: a Greek word in \textit or \textbf keeps its shape and series.
% The transitions to and from every other class carry the tokens that class
% has towards an ordinary letter, so French spacing before « ; » or inside
% « » still holds after a Greek word. No grouping, so no brace or \endgroup
% can come out unbalanced; maths is left alone, since XeTeX inserts nothing
% there. Kept identical in `exercices.sty`.
\makeatletter
\newfontfamily\germainGreekfont{cmun}[
  Extension=.otf, NFSSFamily=germaingreek,
  UprightFont=*rm, ItalicFont=*ti, SlantedFont=*sl,
  BoldFont=*bx, BoldItalicFont=*bi, BoldSlantedFont=*bl]
\newXeTeXintercharclass\germain@greek
\def\germain@greekrange#1#2{%
  \@tempcnta=#1\relax
  \loop
    \XeTeXcharclass\@tempcnta=\germain@greek
    \ifnum\@tempcnta<#2\advance\@tempcnta\@ne\repeat}
\germain@greekrange{"0370}{"03FF}
\germain@greekrange{"1F00}{"1FFF}
\def\germain@greekprev{\rmdefault}
\def\germain@greekin{%
  \edef\germain@greekprev{\f@family}\fontfamily{germaingreek}\selectfont}
\def\germain@greekout{\fontfamily{\germain@greekprev}\selectfont}
\def\germain@greekjoin{%
  \XeTeXinterchartokenstate=\@ne
  \@tempcnta=\z@
  \loop
    \ifnum\@tempcnta=\germain@greek\else
      \XeTeXinterchartoks\germain@greek\@tempcnta=\expandafter{%
        \expandafter\germain@greekout\the\XeTeXinterchartoks\z@\@tempcnta}%
      \XeTeXinterchartoks\@tempcnta\germain@greek=\expandafter{%
        \the\XeTeXinterchartoks\@tempcnta\z@\germain@greekin}%
    \fi
    \ifnum\@tempcnta<4095\advance\@tempcnta\@ne\repeat}
% babel-french sets its own classes, and saves and restores the interchar
% state, each time a language is selected: join again after it does.
\addto\extrasfrench{\germain@greekjoin}
\addto\extrasenglish{\germain@greekjoin}
\AtBeginDocument{\germain@greekjoin}
\makeatother

\usepackage{geometry}
\geometry{a4paper,margin=25mm,marginparwidth=28mm}
\usepackage{marginnote}
\usepackage{xcolor}
\usepackage[normalem]{ulem}
\setlength{\emergencystretch}{3em}
\usepackage{hyperref}
\hypersetup{colorlinks=true,linkcolor=black,urlcolor=black,pdfborder={0 0 0}}

% --- Batch metadata ---------------------------------------------------------
% Taken as they stand by the rendering script, which shows them at the head of
% the reading view. They fix what the file claims to cover. The macro names are
% the parent project's, so that the scripts stay shared; \folder here means the
% volume's slug — `fr-9115` — and \pages counts Gallica views.

\newcommand{\folder}[1]{\gdef\grothFolder{#1}}
\newcommand{\batch}[1]{\gdef\grothBatch{#1}}
\newcommand{\pages}[2]{\gdef\grothFirst{#1}\gdef\grothLast{#2}}
\newcommand{\foldertitle}[1]{\gdef\grothFoldertitle{#1}}

% The holder's own shelfmark, printed as they write it: « Français 9115 ».
\newcommand{\shelfmark}[1]{\gdef\grothShelfmark{#1}}

% The dating, copied verbatim from the holder's catalogue — never inferred
% here. « XIXe siècle » is the BnF's; a pass that dates a leaf from its content
% says so in a \note{}, as its own inference.
\newcommand{\dating}[1]{\gdef\grothDating{#1}}

% The status notice, printed in the title block and repeated as a diagonal
% watermark on every page of the PDF. The manuscripts are in the public
% domain; what this line declares is not a right but a quality — an
% unverified first pass by a machine. The skills write the line; a file
% without it gets no watermark, so leaving it out is visible in the output.
\newcommand{\watermark}[1]{\gdef\grothWatermark{#1}}

\gdef\grothFolder{?}\gdef\grothBatch{}\gdef\grothFirst{?}\gdef\grothLast{?}%
\gdef\grothFoldertitle{}\gdef\grothShelfmark{}\gdef\grothDating{}\gdef\grothWatermark{}

% --- The résumé -------------------------------------------------------------
\newenvironment{resume}
  {\small\setlength{\parskip}{0.45em}}
  {\par\medskip\hrule\medskip}

% --- Keywords ---------------------------------------------------------------
% In English, closing the modernised reading's résumé: the modern vocabulary
% under which to search for what the leaves construct. The manifest extracts
% them to tag the volume — this line is the single source of the tags.
\newcommand{\keywords}[1]{\par\smallskip\noindent
  {\footnotesize\textsc{Keywords} --- \textit{#1}}\par}

% --- The critical apparatus -------------------------------------------------

% The Gallica view number, in the margin. This is the anchor that lets a
% disagreement over a formula be settled by looking at *one* image rather than
% a volume of 750. The site uses it to turn the facsimile as the transcription
% is scrolled. A view is one scanned image: usually one side of a leaf.
\newcommand{\page}[1]{%
  \marginnote{\footnotesize\color{gray!70}\textsf{v.\,#1}}%
  \label{page:#1}\ignorespaces}

% The BnF's pencilled foliation, where the leaf carries it and a pass can read
% it: « 198r ». This is what the literature cites — « ff. 198r–208v » — and
% the one line that turns a citation into a view. Recorded, never inferred:
% a folio a pass cannot read is not written.
\newcommand{\folio}[1]{%
  \marginnote{\footnotesize\color{gray!50}\textsf{f.\,#1}}\ignorespaces}

% The modernised reading's coarser counterpart to \page{}: which views a
% section is drawn from.
\newcommand{\pagerange}[2]{%
  \marginnote{\footnotesize\color{gray!70}\textsf{v.\,#1--#2}}%
  \ignorespaces}

% Illegible: never guessed.
\newcommand{\ill}{\textcolor{red!60!black}{[\,\ldots\,]}}

% A doubtful reading: the word is given, and flagged as uncertain.
\newcommand{\uncertain}[1]{\uline{#1}}

% An editorial addition — a missing letter, an implied word.
\newcommand{\add}[1]{\textcolor{blue!50!black}{[#1]}}

% Struck out by the writer. What was crossed out is part of the document.
\newcommand{\struck}[1]{\sout{#1}}

% The transcriber's note, on the state of the page or the mathematics.
\newcommand{\note}[1]{\par\smallskip{\small\color{gray!60!black}\textit{[#1]}}\par\smallskip}

% A marginal note of hers, distinct from the body of the text.
\newcommand{\marginal}[1]{\par\smallskip\noindent\hspace{1em}%
  {\small\color{gray!60!black}#1}\par\smallskip}

% --- Title ------------------------------------------------------------------
% The title block is French, like the documents themselves.

\AddToHook{shipout/background}{%
  \ifx\grothWatermark\empty\else
    \begin{tikzpicture}[remember picture, overlay]
      \node[rotate=52, scale=3.4, align=center, text=gray!18,
            font=\sffamily\bfseries]
        at (current page.center) {\grothWatermark};
    \end{tikzpicture}%
  \fi}

\AtBeginDocument{%
  \begin{center}
    {\large\bfseries Manuscrits de mathématiciens (Gallica) ---
      \ifx\grothShelfmark\empty\grothFolder\else\grothShelfmark\fi}\\[2pt]
    {\itshape\grothFoldertitle}\\[4pt]
    {\small vues \grothFirst--\grothLast
      \ifx\grothBatch\empty\else\ (lot \grothBatch)\fi}\\[2pt]
    \ifx\grothDating\empty\else
      {\small\color{gray!60!black}Datation du catalogue : \grothDating}\\[2pt]
    \fi
    \ifx\grothWatermark\empty\else
      {\small\bfseries\color{red!45!gray}\grothWatermark}\\[2pt]
    \fi
    {\footnotesize\color{gray!60!black}Transcription établie d'après les images IIIF de Gallica.
     Source gallica.bnf.fr / Bibliothèque nationale de France.\\
     Les lectures incertaines sont soulignées, les illisibles notées [\,\ldots\,],
     les ajouts entre crochets ; \textsf{v.} numérote les vues de Gallica,
     \textsf{f.} la foliotation au crayon de la BnF.}
  \end{center}
  \vspace{1em}}

\endinput


\folder{fr-9118}
\shelfmark{Français 9118}
\pages{1}{103}
\dating{1801-1900}
\watermark{Lecture automatique\\interprétation non vérifiée}
\foldertitle{Correspondance de Sophie GERMAIN avec les mathématiciens et savants Cauchy, Delambre, Fourier, Gauss, Le Gendre, d'Ansse de Villoison, etc. — lecture modernisée du volume entier}

\begin{document}

\section*{Résumé}

\begin{resume}
Ce volume n'est pas un cahier de mathématiques. C'est le dossier extérieur
d'une carrière : cent trois vues où l'on trouve les lettres qu'elle a reçues —
Cauchy, Delambre, Fourier, Legendre, Navier, Poisson, Gauss, d'Ansse de
Villoison, et quelques autres —, trois brouillons des siennes, la copie de la
note par laquelle Lagrange a jugé son mémoire de 1811, un cours de mécanique
élémentaire écrit pour elle par une main qui ne se nomme pas, une plaquette
imprimée, trois récépissés datés de l'Institut, et une douzaine de feuillets
détachés de sa main. Les mathématiques qu'on y lit sont donc, pour la plus
grande part, celles que d'autres lui écrivent. C'est ce qui fait leur prix :
elles disent ce qu'on lui a objecté, et ce qu'elle a répondu.

Deux questions occupent presque tout. La première est physique : quelle
équation gouverne les vibrations d'une plaque élastique ? En 1811 l'Institut
met la question au concours ; elle est la seule à concourir ; son mémoire est
jugé faux, et la note de Lagrange qu'on trouve ici, copiée, donne à sa place
l'équation qui est encore la nôtre. La seconde est arithmétique : les
\emph{Disquisitiones arithmeticae} de Gauss, qu'elle lit dès leur parution,
lui fournissent la matière de deux longues lettres — écrites sous le nom
d'emprunt « Le Blanc » — dont les brouillons sont ici, avec les réponses de
Gauss.

De ce qu'elle cherchait, elle a laissé dans ce volume son propre compte rendu.
Un brouillon non daté, « ce 18 juillet », énumère les trois choses de son
« petit écrit » qui lui semblent « de quelqu'importance » : la définition
exacte de la question — quelles expériences en relèvent et lesquelles n'en
relèvent pas ; les couches dont l'épaisseur d'une plaque est composée ; et la
courbure moyenne, qu'elle compare à une température moyenne. Cette dernière
comparaison, qu'elle avance comme une commodité de calcul, est plus juste
qu'elle ne pouvait le savoir : la somme des deux courbures principales
$1/r + 1/r'$ est, pour une surface presque plane $z(x,y)$, le laplacien
$\Delta z$ — l'opérateur même de la théorie de la chaleur qu'elle invoque par
analogie. L'énergie de flexion $\iint (\Delta z)^2$ conduit alors, en une
ligne de calcul des variations, à l'équation
$\partial^2 z/\partial t^2 + k^2\,\Delta^2 z = 0$ que Lagrange écrit dans sa
note et que Legendre lui transmet. Son hypothèse physique était la bonne ;
c'est son analyse qui ne l'était pas, et Poisson le lui écrit en toutes
lettres en 1816.

Du côté de l'arithmétique, les brouillons à Gauss portent trois choses. Une
généralisation du théorème de l'article 357 des \emph{Disquisitiones} — que
$4(x^n-1)/(x-1)$ est de la forme $y^2 \pm n z^2$ — qu'elle étend à
$4(x^{n^s}-1)/(x-1)$, et qui est vraie : elle se déduit en une ligne de la
composition des formes. Une conjecture sur les formes quadratiques à $n$
variables, que le déterminant de l'adjointe vaut $D^{n-1}$, qui est vraie
aussi et qui est aujourd'hui une identité de première année. Et une annonce
qu'elle raye puis récrit telle quelle : qu'elle croit avoir démontré
l'impossibilité de $x^n + y^n = z^n$ pour $n = p - 1$, $p$ premier de la forme
$8k+7$. La démonstration était jointe à la lettre ; elle n'est pas dans le
volume, et rien ici ne permet de la juger. Trois feuillets détachés, non
datés, y ajoutent une objection à Gauss sur la composition des formes binaires,
et deux remarques sur $p^4+4$ et $p^4+q^4$ — la première étant, sous une autre
démonstration que la sienne, ce qu'on appelle aujourd'hui l'identité de Sophie
Germain.

Où cela mène. Du côté des plaques : à l'opérateur biharmonique $\Delta^2$, à
la théorie des plaques minces de Kirchhoff et Love, et aux conditions de bord
qui ont occupé le reste du siècle. Du côté de la lame vibrante : aux modes
propres de la poutre d'Euler-Bernoulli, dont les nœuds ne sont pas également
espacés — ce qu'elle oppose à Euler sur un feuillet, et Legendre sur trois.
Du côté de l'arithmétique : aux polynômes cyclotomiques et aux sommes de
Gauss, à la composition des formes quadratiques binaires et au groupe des
classes, et — pour l'équation de Fermat — à une histoire qui ne s'achève
qu'en 1995.

\keywords{Kirchhoff–Love plate theory, biharmonic operator, mean curvature,
Euler–Bernoulli beam vibrations, nodal lines, cyclotomic polynomials,
composition of binary quadratic forms, Sophie Germain's identity, Fermat's
last theorem, central binomial coefficient}
\end{resume}

\section*{Ce que le volume contient}

\pagerange{1}{103}
Cent trois vues, dont les quatre premières et les trois dernières sont la
reliure et les gardes.\footnote{Les vues 101 à 103 — verso blanc du dernier
feuillet, garde et contre-garde, plat inférieur — ne portent aucun texte ;
c'est ce que constate le lot 6 de la transcription, qui ne transcrit rien et
le dit. Le texte du volume s'arrête à la vue 100.} Entre les deux,
quatre-vingt-treize feuillets numérotés au crayon par la Bibliothèque, qui se
rangent en six espèces.

\begin{itemize}
  \item \emph{Les lettres reçues}, qui font la moitié du volume : deux de
    Cauchy (1821, 1826), deux de Delambre (1816, 1821), sept de Fourier (une
    officielle de 1823, une de 1824, une de 1826, et quatre autographes non
    datées), huit de Legendre (de 1811 à 1819, plus celle de 1821), une de
    Navier (1821), une de Poisson (1816), trois de Gauss (1805, 1805, 1808),
    une du général Pernety avec la lettre d'un officier qu'il avait envoyé
    chez Gauss (1806), trois de d'Ansse de Villoison (1802), un billet d'un
    libraire nommé Bernard, et une invitation à dîner non signée. Presque
    toutes ont conservé leur feuillet d'adresse, plié et cacheté, et ces
    adresses sont ici le seul registre de ses domiciles successifs.
  \item \emph{Deux notes mathématiques de Legendre} sans en-tête ni date,
    sur un problème d'Euler (ff. 42r–43r et ff. 58–59, vues 47 à 48 et 63 à
    64) ; et \emph{une note signée G. Libri}, datée 1824 (f. 60, vue 65).
  \item \emph{Trois brouillons de sa main} : une lettre « ce 18 juillet » à un
    correspondant que le feuillet ne nomme pas (ff. 1r–2r, vues 5 et 6), et
    deux lettres à Gauss signées « Le Blanc » (ff. 38r–41v, vues 43 à 47).
  \item \emph{Une copie}, d'une main qui n'est pas la sienne, de la « Note
    communiquée aux commissaires pour le prix de la surface élastique
    (décembre 1811) », avec la mention « cette note est de Mr la grange »
    (f. 34r, vue 39).
  \item \emph{Un cours}, « Réponse aux questions de Mlle \ldots », d'une
    troisième main : le levier, l'infini, l'hyperbole, $0/0$, et deux
    indications de librairie (ff. 35r–36v, vues 40 à 42).
  \item \emph{Des pièces imprimées et administratives} : une plaquette de huit
    pages, les vers latins de Villoison pour l'anniversaire de Lalande, où une
    note en bas de page la nomme (ff. 74–78, vues 81 à 85) ; et trois
    récépissés de l'Institut, datés du 21 septembre 1811, du 21 septembre 1813
    et du 8 novembre 1821 (ff. 78, 80, 82, vues 85 à 89).
  \item \emph{Des feuillets détachés de sa main}, à la fin : deux rédactions
    d'un problème de centre de gravité, trois notes de lecture
    d'arithmétique, trois feuillets de physique d'une écriture plus ronde, et
    un feuillet contre Euler sur les lignes nodales (ff. 84–93, vues 91 à 99).
\end{itemize}

Le classement n'est pas chronologique et n'est pas thématique : les lettres de
Legendre de 1811, 1813 et 1819 se suivent, mais celle de 1821 est vingt vues
plus loin ; les feuillets d'arithmétique et ceux de physique alternent sans
raison lisible. Ce qui suit est rangé par matière, non par vues.

\section*{Les surfaces élastiques : ce qu'elle cherchait, dans ses termes}

\pagerange{5}{6}
Le brouillon « ce 18 juillet » est la seule page du volume où elle explique
son propre projet. Elle vient d'imprimer un « petit écrit » — ses « nouvelles
remarques » — sur le conseil du correspondant, et lui en signale trois
points.\footnote{Le feuillet ne nomme pas le destinataire. Il parle des
expériences de Savart, mentionne Ampère comme intermédiaire, et demande où
trouver ce que le correspondant a publié « sur le cas général du mouvement des
corps élastiques ». Le mot « nouvelles remarques » est celui du titre que
Fourier accuse réception le 24 juillet 1826 — les \emph{Remarques sur la
nature, les bornes et l'étendue de la question des surfaces élastiques} — et
c'est Cauchy qui fut alors désigné pour en faire un rapport verbal. Cela ne
suffit pas à identifier le destinataire, et la présente lecture ne
l'identifie pas.}

Le premier est la délimitation de la question. Elle tient à dire que les
expériences de Savart sont « étrangères à cette question » et ne pourroient
s'expliquer que par une théorie plus étendue : ce qu'elle traite est un cas
particulier du mouvement des solides élastiques, non la totalité du
phénomène. C'est, en langage d'aujourd'hui, la déclaration d'un domaine de
validité — la théorie linéaire des petits déplacements d'une plaque mince —
et l'aveu que ce qui sort de ce domaine sort de sa théorie.

Le deuxième est l'épaisseur. « Si dans les expériences dont je parle on
parvenoit à séparer les différentes couches dont on peut concevoir que
l'épaisseur soit composée, chacune d'elles présenteroit des figures
particulières. » C'est la question de la variation des grandeurs à travers
l'épaisseur, à laquelle elle consacrera un mémoire entier, le \emph{Mémoire
sur l'emploi de l'épaisseur dans la théorie des surfaces élastiques} lu le 9
mars 1824.\footnote{La lettre officielle de Fourier du 15 mars 1824 (f. 16r,
vue 20) en donne le titre, la date de lecture et la commission : Laplace,
Prony, Poisson. Le billet autographe « vendredi matin » (ff. 24r–25r, vues 28
et 29) rapporte la même séance et la même commission ; la transcription en
infère qu'il serait du 12 mars 1824, et le dit inférence.} La cinématique
moderne des plaques minces — les fibres normales restent droites et normales
— est précisément une hypothèse sur ces couches ; elle est due à Kirchhoff
(1850) et à Love, et l'anachronisme doit être marqué : le feuillet ne l'a pas.

Le troisième est la courbure moyenne, et c'est le seul des trois qui soit une
idée mathématique. Elle écrit que « l'état réel des points dont la position
est également éloignée de ceux auxquels appartiennent les manières d'être
extrêmes donne l'état moyen du système », et prend son exemple dans un anneau
chauffé : la température des points également éloignés du maximum et du
minimum est la température moyenne de la pièce. De même, dit-elle, il faut des
forces égales pour changer une courbure moyenne donnée, « ensorte que la
courbure de la sphère est toujours comparable à celle d'une surface de figure
quelconque » — et cette remarque « m'a servi à donner une forme assez simple à
l'équation générale des surfaces ».

En termes actuels : la quantité $1/r + 1/r'$, somme des deux courbures
principales, est deux fois la courbure moyenne $H$, c'est-à-dire la trace de
l'opérateur de forme. C'est bien une moyenne, au sens exact : la courbure
normale, prise dans une direction faisant l'angle $\theta$ avec la première
direction principale, vaut $k_1\cos^2\theta + k_2\sin^2\theta$, dont la
moyenne sur toutes les directions est $(k_1+k_2)/2 = H$. C'est aussi une
grandeur intrinsèquement associée au point, indépendante du repère choisi,
et la sphère est la surface où elle est constante. Et pour un graphe
$z(x,y)$ presque plan,
\[
  \frac{1}{r} + \frac{1}{r'} \;=\; \Delta z + O\!\left(|\nabla z|^2\right),
  \qquad \Delta = \frac{\partial^2}{\partial x^2} + \frac{\partial^2}{\partial y^2},
\]
de sorte que sa mesure de la courbure est, au premier ordre, le laplacien —
l'opérateur même qui gouverne la propagation de la chaleur dans l'anneau
qu'elle prend pour terme de comparaison.\footnote{L'analogie est la sienne ;
l'identification de $1/r+1/r'$ au laplacien linéarisé est de la présente
lecture. Le mot du feuillet est « l'armille », que la transcription souligne
comme une lecture douteuse. Rien sur la page ne dit qu'elle ait vu dans les
deux cas le même opérateur : elle dit seulement que la même manière de
raisonner s'y applique.}

\section*{L'équation de la plaque : la note de Lagrange et la lettre de Legendre du 4 décembre 1811}

\pagerange{39}{54}
C'est le nœud du volume, et il s'y trouve deux fois : dans la copie de la note
de Lagrange (f. 34r, vue 39) et dans la lettre où Legendre la lui rapporte
(ff. 48r–49r, vues 53 et 54). Les deux textes donnent la même équation.

\subsection*{Ce que Lagrange écrit}

La note commence par un refus : « L'équation fondamentale pour le mouvement de
la surface vibrante ne me paraît pas exacte, et la manière dont on cherche à
la déduire de celle d'une lame élastique en passant d'une ligne à une surface
ne paraît peu juste. » Elle donne ensuite ce à quoi se réduit, pour $z$ très
petit, l'équation du mémoire :
\[
  \frac{\partial^2 z}{\partial t^2}
  + C\left(\frac{\partial^6 z}{\partial x^4\,\partial y^2}
         + \frac{\partial^6 z}{\partial y^4\,\partial x^2}\right),
\]
puis, en adoptant $1/r + 1/r'$ comme mesure de la courbure à laquelle
l'élasticité est supposée proportionnelle, celle qu'il trouve lui-même :
\[
  \frac{\partial^2 z}{\partial t^2}
  + k^2\left(\frac{\partial^4 z}{\partial x^4}
           + 2\,\frac{\partial^4 z}{\partial x^2\,\partial y^2}
           + \frac{\partial^4 z}{\partial y^4}\right) = 0,
\]
« qui est bien différente de la précédente ».\footnote{Deux écarts avec le
feuillet. Le coefficient de la première équation y est lu « gEbc », dont la
première lettre pourrait aussi être un 9 ; il est écrit $C$ ici parce qu'aucune
des deux lectures ne se laisse interpréter, et rien du raisonnement n'en
dépend. Et la première équation, sur le feuillet, n'est pas une équation : elle
s'achève sur un point-virgule, sans second membre. On la lit comme
« $\dots = 0$ » par symétrie avec la seconde, mais c'est une supposition, et la
transcription ne la fait pas.}

La seconde équation est, en écriture d'aujourd'hui,
\[
  \frac{\partial^2 z}{\partial t^2} + k^2\,\Delta^2 z = 0,
\]
où $\Delta^2 = \partial^4/\partial x^4 + 2\,\partial^4/\partial x^2\partial y^2
+ \partial^4/\partial y^4$ est l'opérateur biharmonique.\footnote{Lagrange
écrit le $d$ rond, Legendre le $d$ droit ; les deux notent la même dérivée
partielle et sont rendus ici par $\partial$. Le nom d'« opérateur
biharmonique » et la notation $\Delta^2$ sont postérieurs aux feuillets.} La
forme développée n'est pas seulement plus courte : elle rend visible que
l'opérateur est le carré du laplacien, donc invariant par rotation du plan —
ce que la somme de deux dérivées sixièmes croisées n'est pas.

Le chemin qui y mène est celui que l'hypothèse de Germain impose, et il tient
en quelques lignes de calcul des variations. Si l'élasticité est
proportionnelle à
$1/r+1/r'$, l'énergie de flexion est, à petit déplacement,
\[
  E[z] = \frac{1}{2}\iint \left(\Delta z\right)^2 dx\,dy,
\]
dont la variation première est $\iint (\Delta^2 z)\,\delta z\,dx\,dy$ plus des
termes de bord ; le principe de d'Alembert donne alors
$\partial_t^2 z = -k^2\Delta^2 z$.\footnote{Cette dérivation est de la présente
lecture ; ni la note de Lagrange ni la lettre de Legendre ne l'écrivent — la
note dit seulement « je trouve », et Legendre ajoute « Je n'ai point vérifié ce
calcul ». Elle est donnée ici parce qu'elle montre que l'équation de Lagrange
suit de l'hypothèse de Germain, et que le désaccord de 1811 porte donc sur
l'analyse et non sur la physique.} L'équation moderne de la plaque mince,
$\rho h\,\partial_t^2 w + D\,\Delta^2 w = 0$, est la même à la constante près.

\subsection*{La vérification par la lame}

Legendre, qui rapporte la conclusion sans l'avoir vérifiée, donne aussitôt la
raison qui l'en persuade : « en supposant la surface vibrante réduite à une
lame d'une largeur constante, ce qui peut s'exprimer en faisant $dz/dy = 0$,
on retombe sur l'équation
$\partial^2 z/\partial t^2 + k^2\,\partial^4 z/\partial x^4 = 0$ », celle
d'Euler pour les lames élastiques vibrantes ; « Votre équation ne donneroit pas
ce résultat. »

Il a raison, et l'argument est le bon : c'est un test de cohérence par passage
à la limite d'une dimension. Si $z$ ne dépend pas de $y$, alors
$\Delta^2 z = \partial^4 z/\partial x^4$ et l'équation de la plaque devient
celle de la poutre d'Euler-Bernoulli ; tandis que la première équation, où
chaque terme porte au moins une dérivée en $y$, se réduit à
$\partial_t^2 z = 0$, c'est-à-dire à rien.\footnote{La condition qu'il écrit,
$\partial z/\partial y = 0$, ne suffit pas telle quelle : elle n'annule que la
dérivée première, et les termes en $\partial^4/\partial x^2\partial y^2$
demandent que les dérivées supérieures en $y$ s'annulent aussi. Ce qu'il faut
est que $z$ ne dépende pas de $y$, ce qui est manifestement ce qu'il veut dire
en parlant d'« une lame d'une largeur constante ». La correction est de la
présente lecture.}

\section*{Où l'analyse a manqué : les intégrales doubles et les « quatre points »}

\pagerange{55}{56}
La lettre du 4 décembre 1813 est la plus dure du volume, et la plus précise.
Legendre y dit trois choses.

D'abord que Lagrange a eu raison de considérer, dans la courbe élastique, deux
éléments consécutifs et de mesurer l'élasticité par l'angle qu'ils
comprennent. C'est exact : pour une courbe plane, l'angle entre deux éléments
consécutifs divisé par la longueur d'arc \emph{est} la courbure, et l'énergie
$\int \kappa^2\,ds$ est celle de l'élastique.

Ensuite qu'« on n'a pas d'Elemens analogues dans les surfaces ». Il précise :
un élément de surface a pour projection $dx\,dy$, le suivant
$(dx+ddx)(dy+ddy)$, « ces projections sont deux carrés séparés » ; et
« l'idée des plans ne s'accommode pas avec ces projections, parcequ'un plan ne
passe pas par quatre points ». En termes d'aujourd'hui : la courbure d'une
surface en un point n'est pas un nombre mais une forme quadratique — la
seconde forme fondamentale — et il n'existe pas d'« angle entre deux éléments
consécutifs » qui la représente, parce qu'un quadrilatère infinitésimal tracé
sur une surface n'est en général pas plan. Le défaut de planéité est
précisément d'ordre deux, c'est-à-dire du même ordre que ce qu'on veut
mesurer ; passer de la ligne à la surface par transport de l'angle ne peut
donc pas marcher. La bonne quantité scalaire est la trace de cette forme,
c'est-à-dire $1/r+1/r'$ — ce que son hypothèse disait déjà — et il faut la
faire entrer dans une intégrale double avant de varier, non varier la ligne
puis intégrer.

Enfin que c'est « dans les doubles intégrales » qu'elle s'est égarée, et que
la méthode à suivre était celle de Lagrange, première édition, page
148.\footnote{L'identification de l'ouvrage — la \emph{Mécanique analytique} —
est celle de la transcription, non du feuillet.} Il ajoute, et c'est la
concession qui compte : « Au reste ces choses sont sujettes à des difficultés
particulieres qui n'ont pas encore été bien éclaircies, et il y auroit des
objections à faire contre l'analyse de l'article même que je cite. »

Le partage entre l'hypothèse et le calcul, que la présente lecture a tiré de
la note de Lagrange, est dit textuellement par Poisson dans sa lettre du 15
janvier 1816 (ff. 67–68, vues 72 à 74) : « Le reproche que la Commission lui a
fait porte moins sur l'hypothèse dont vous êtes partie que sur la manière dont
vous avez appliqué le calcul à cette hypothèse. » Il ajoute que son propre
résultat et le sien ne s'accordent que « dans le seul cas où la surface
s'écarte infiniment peu d'un plan ».\footnote{« le mien » est une lecture
douteuse de la transcription, qui propose aussi « l'ancien ». Le sens du reste
de la phrase — la comparaison de « mes résultats aux vôtres » deux lignes plus
bas — favorise « le mien », mais la présente lecture ne tranche pas.} C'est
exactement la situation moderne : les deux voies, la sienne par la courbure et
celle de Poisson par les forces moléculaires, donnent la même équation
linéaire de la plaque et diffèrent hors du régime linéaire, où l'équation
n'est plus $\partial_t^2 z + k^2\Delta^2 z = 0$.

Deux lettres plus tard, le désaccord a changé d'objet. Fourier, dans le billet
autographe « vendredi matin » (ff. 24r–25r, vues 28 et 29), écrit : « il m'a
paru que vous mettez dans tout son jour l'insuffisance de l'hypothèse
théorique dont on a voulu déduire l'équation du 4\textsuperscript{e} ordre que
vous avez trouvée ». L'équation du quatrième ordre lui est ici attribuée ; ce
qui est en cause n'est plus elle, mais l'hypothèse dont on prétend la
déduire.\footnote{La lettre de Legendre du 4 décembre 1813 va dans le même
sens : « Il paroit reconnu cependant que votre Equation est réellement celle de
la surface vibrante ; en mettant l'analyse à part le reste peut etre bon. » Ni
l'une ni l'autre lettre ne dit de quelle équation il s'agit — celle du mémoire
de 1811, ou celle que la correction de Lagrange a substituée ; la présente
lecture ne le décide pas.}

\section*{Euler et la lame vibrante : trois notes de Legendre}

\pagerange{57}{64}
Trois pièces de Legendre, dont aucune ne porte de date lisible et sûre,
discutent le même texte d'Euler sur les vibrations d'une lame élastique, cité
par pages (152 à 157) et par paragraphes (§ 45, § 47).\footnote{Aucun des
feuillets ne nomme l'ouvrage. La transcription du lot 3 propose le traité
d'Euler sur les courbes élastiques et le donne pour son identification ; la
présente lecture s'en tient à « le texte d'Euler que les lettres citent par
pages ».} Elles ne s'accordent pas entre elles, et l'ordre dans lequel il faut
les lire n'est pas établi — voir la dernière section.

\subsection*{La racine parasite $\sin\frac{1}{2}\omega = 0$}

La note intitulée « Pag. 155 » (ff. 58–59, vues 63 et 64) part de l'équation
générale d'Euler, que Legendre récrit\footnote{Legendre surmonte d'un trait
l'exposant négatif — son $\bar{e}^{x}$ est $e^{-x}$ — et se sert du même trait
comme d'une parenthèse ; il écrit « cot. » avec le point. La transcription
conserve ces signes ; ils sont traduits ici sans autre avis.}
\[
  0 = 2 - 2e^{2\omega}
    - \left(e^{\lambda\omega} - e^{(2-\lambda)\omega}\right)
      \left(e^{\lambda\omega} - e^{-\lambda\omega}\right)
      \left(\cot \lambda\omega + \cot (1-\lambda)\omega\right).
\]
Sa remarque est que $\sin\frac{1}{2}\omega = 0$ n'en est pas une racine :
« elle vient d'un facteur qui a été introduit par la Multiplication et qui est
étranger à la solution du problème ». Il a raison, et cela se voit
directement. Pour $\lambda = \frac12$, le troisième facteur est
$2\cot\frac{\omega}{2}$, qui devient infini quand $\sin\frac{\omega}{2}$
s'annule, tandis que les deux autres restent finis et non nuls ; le produit
diverge et l'égalité ne peut avoir lieu. Legendre ajoute la seule exception :
elle « le seroit seulement par la supposition $\frac12\omega = 0$, ou
$\frac12\omega =$ à un infiniment petit » — et en effet, quand $\omega\to 0$,
$e^{\omega/2}-e^{-\omega/2}\sim\omega$ compense le $\cot\frac\omega2\sim 2/\omega$
et tout tend vers $0$, mais ce cas est celui du repos.

Il donne ensuite la raison qui rend cette racine inadmissible même comme
solution formelle : elle « donne des valeurs infinies pour les coefficiens
$\gamma$ et $\gamma'$, $\delta'$ ». Et il reproche à Euler d'avoir écrit
« Multiplicemus omnes Coefficientes par $\sin\frac12\omega$ » : « on peut bien
multiplier l'ordonnée d'une Courbe par une Constante \ldots{} mais on ne peut
pas multiplier par zéro ». C'est exact, et c'est le diagnostic moderne d'une
racine parasite : une équation obtenue en multipliant par une quantité qui
peut s'annuler acquiert les zéros de cette quantité, qui ne sont pas des
solutions du problème.

Il porte le même reproche à Germain : « Lorsque Mlle Sophie a
voulu considérer le cas général, elle est ce me semble tombée dans la même
erreur qu'Euler, en faisant $\sin\lambda\omega = 0$. »

\subsection*{La résolution générale, et le cas $\lambda = 1/3$}

La même note donne alors une méthode pour résoudre l'équation d'Euler sans
tâtonnement, et c'est la plus jolie page de mathématiques du volume. Passé la
première racine, dit Legendre, $e^{\omega}$ est si grand qu'on peut négliger
$e^{-\omega}$ devant $e^{\omega}$ et $e^{-\lambda\omega}$ devant
$e^{\lambda\omega}$ ; l'équation se réduit à
\[
  0 = -2e^{2\omega} + e^{2\omega}\left(\cot\lambda\omega + \cot(1-\lambda)\omega\right),
  \qquad\text{soit}\qquad
  \cot\lambda\omega + \cot(1-\lambda)\omega = 2 .
\]
C'est un développement asymptotique, mené comme on le mènerait aujourd'hui :
$e^{\lambda\omega}-e^{(2-\lambda)\omega} \sim -e^{(2-\lambda)\omega}$ et
$e^{\lambda\omega}-e^{-\lambda\omega} \sim e^{\lambda\omega}$, dont le produit
est $\sim -e^{2\omega}$, et $2-2e^{2\omega}\sim -2e^{2\omega}$.

Posant $x = \cot\lambda\omega$ et $y = \cot(1-\lambda)\omega$, il observe que
pour chaque valeur rationnelle de $\lambda$ une relation algébrique lie $x$ et
$y$, laquelle, jointe à $x+y=2$, donne un nombre fini de racines
$\lambda\omega = \alpha, \beta, \ldots$, d'où les suites
$\lambda\omega = \alpha + k\pi$, $\lambda\omega = \beta + k\pi$, \ldots{} Il
traite $\lambda = \frac13$ : avec $x = \cot\frac{\omega}{3}$, la formule de
duplication donne $\cot\frac{2\omega}{3} = (x^2-1)/2x$, d'où
\[
  x + \frac{x^2-1}{2x} = 2, \qquad 3x^2 - 4x - 1 = 0, \qquad
  x = \frac{2 \pm \sqrt{7}}{3},
\]
et $\omega$ décrit les deux suites $3\alpha + 3k\pi$ et $3\beta + 3k\pi$, où
$\cot\alpha = (2+\sqrt7)/3$ et $\cot\beta = (2-\sqrt7)/3$. Le calcul est juste
d'un bout à l'autre.\footnote{Le feuillet écrit l'équation « $3x^2 - 1 = 4x$ »,
qui est la même. Numériquement, $\alpha \approx 0{,}573$ et
$\beta \approx 1{,}781$ radians, d'où $\omega \approx 1{,}72$ et
$\omega \approx 5{,}34$ pour $k = 0$.} Legendre ferme d'ailleurs lui-même la
seule faiblesse de la méthode : « Dans l'application il faudroit rechercher
plus exactement les deux premiers termes $3\alpha$, $3\beta$, mais les autres
seront toujours suffisamment approchés » — c'est-à-dire que l'approximation,
qui suppose $e^{\omega}$ grand, est médiocre pour les premières racines et
s'améliore avec $k$. C'est exactement ce qu'on dirait aujourd'hui d'un
développement pour les grands modes.

\subsection*{$\gamma = \gamma'$ : les deux portions ne font qu'une courbe}

La lettre du 19 janvier (ff. 52r–53r, vues 57 et 58) revient sur une erreur de
signe d'Euler. Euler trouve $\gamma' = -\gamma$ ; Legendre soutient que les
équations III et IV de la page 152 donnent $\gamma = \gamma'$, et en tire la
conséquence : l'équation de la portion de courbe $LF$ n'est pas
$y = -C\sin(\zeta + \cdots)$ mais
\[
  y = C \sin\!\left(\zeta + t\,\frac{\omega c \sqrt{2gb}}{a^2}\right)\sin u\omega,
\]
« absolument comme celle de la portion $EL$, c'est-à-dire que ces deux
portions ne font qu'une seule et même courbe désignée par la même
équation ».\footnote{Le numérateur $\omega c\sqrt{2gb}$ est d'une lecture
douteuse, et le feuillet écrit $aa$ pour $a^2$. La transcription note aussi
que « $LF$ » est écrit par-dessus un « $EF$ ». Rien de ce qui suit ne dépend de
la valeur de ce coefficient, qui n'est qu'une pulsation.}

L'objection qu'on pourrait faire — que l'équation d'Euler, avec son signe
moins, indique bien les ordonnées négatives au-delà du nœud — il la lève
lui-même : « puisque $\sin u\omega$ est zéro lorsqu'on fait $u = \frac12$, les
deux suppositions $u > \frac12$, $u < \frac12$ donneront deux résultats de
signes contraires pour $\sin u\omega$, desorte qu'avant et après le point $L$,
les ordonnées seront de signes différens ».

C'est le raisonnement juste, et c'est celui qu'on ferait aujourd'hui : un mode
propre est \emph{une} fonction, analytique sur toute la longueur ; un nœud est
un zéro de cette fonction, et le changement de signe des ordonnées de part et
d'autre est celui du facteur qui s'annule, non un changement de la loi. Écrire
deux équations pour les deux arcs, avec des constantes de signes opposés, c'est
compter deux fois le même changement de signe.

\section*{Sa propre objection à Euler sur les lignes nodales}

\pagerange{99}{99}
Un feuillet détaché de sa main (f. 93, vue 99) attaque le même auteur sur le
même point, et par l'autre bout. Euler établit, dit-elle, que les angles que
les deux portions font avec l'axe au nœud $L$ doivent être égaux, et que les
rayons osculateurs y sont les mêmes ; elle objecte que pour des extrémités
libres ou fixes cela n'est pas valable dès que le paramètre diffère de
$\frac12$, « puisque j'ai trouvé qu'alors les courbures prises à droite et à
gauche du point $L$ étoient différentes entr'elles, ce qui résulte de ce
qu'aucune ligne nodale située ailleurs qu'au milieu de la lame n'est à
distances égales des deux plus voisines ». Elle ajoute qu'il lui semble que
$dy/ds$ et $d^2y/ds^2$ — « les tangentes et les rayons osculateurs » —
« doivent être de signes différens pour les deux portions ».\footnote{Le
feuillet écrit la dérivée seconde $ddy/ds^2$, avec le $d$ droit ; elle est
rendue ici par $d^2y/ds^2$. Le renvoi entre parenthèses, « § 45 de la
P. », est d'une lecture douteuse dans la transcription, qui ne parvient
ni à lire le mot abrégé ni à assurer la lettre.}

Il faut séparer deux affirmations.

La première — que les nœuds ne sont pas également espacés — est vraie, et
c'est le fait qui distingue la poutre de la corde. Les modes d'une corde sont
des sinus, dont les zéros sont en progression arithmétique ; ceux d'une poutre
sont des combinaisons de $\sin$, $\cos$, $\sinh$ et $\cosh$, dont les zéros ne
le sont pas — ils sont seulement disposés symétriquement de part et d'autre du
milieu, lorsque les conditions aux deux bouts sont les mêmes. Il en résulte, comme elle
le dit, que les deux arcs qu'un nœud sépare n'ont pas la même longueur et ne
sont donc pas superposables : la figure d'Euler, où les deux portions sont
image l'une de l'autre, ne convient qu'au cas symétrique.

La seconde — que la courbure prise à droite et à gauche de $L$ diffère — ne
tient pas telle qu'elle est écrite. Un mode est une seule fonction, de classe
$C^{\infty}$ ; en un nœud intérieur, la tangente et le rayon osculateur ont
une valeur et une seule.\footnote{C'est exactement la conclusion de la lettre
de Legendre du 19 janvier lue ci-dessus, à laquelle elle s'oppose ici sans le
savoir ou sans le dire. Que le volume tienne les deux feuillets est ce qui
permet de les mettre en regard ; leur ordre chronologique n'est pas établi.}
Ce qui est vrai, et qui est sans doute ce qu'elle vise, c'est que les deux
\emph{arcs} sont de courbures opposées : $y$ change de signe en $L$, donc
$d^2y/ds^2$ aussi de part et d'autre, et la concavité est tournée dans un sens
sur un arc et dans l'autre sur le suivant. Sa seconde phrase est donc juste
comme énoncé sur les deux portions, et fausse comme énoncé au point $L$
lui-même. On notera que $dy/ds$, lui, ne change pas de signe au nœud : c'est
la pente, qui y est commune aux deux arcs.

\section*{L'arithmétique : les deux brouillons à Gauss}

\pagerange{43}{47}
Deux brouillons de sa main, sans date, signés « Le Blanc » pour le premier
(ff. 38r–41v, vues 43 à 47). Ils sont fortement corrigés, avec des débuts de
phrase annulés et des additions interlinéaires, et la transcription en
signale plusieurs passages incertains ; ce sont des brouillons, non des
lettres.\footnote{Aucun des deux feuillets ne porte de date. La transcription
infère que le premier est le brouillon de la lettre à laquelle Gauss répond le
16 juin 1805, et que le mémoire dont le second remercie est celui de 1799
annoncé par Gauss le 20 août 1805 ; elle donne les deux pour ses inférences.
Les deux lettres de Gauss sont dans le volume (ff. 28r–30r, vues 32 à 35), et
l'adresse « chez Mr. Sylvestre de Sacy, Rue haute feuille N° 11 » qu'elle
indique au bas du premier brouillon est celle que porte la lettre de Gauss.}

\subsection*{Le théorème de l'article 357 et la généralisation qu'elle propose}

Elle part du « beau théorème contenu dans l'équation
$4\frac{(x^n-1)}{x-1} = y^2 \pm n z^2$ » et propose de le généraliser ainsi :
\[
  4\,\frac{x^{n^s}-1}{x-1} = y^2 \pm n z^2,
\]
« $n$ étant toujours un nombre premier et $s$ un nombre quelconque ». Elle
ajoute qu'elle joint à sa lettre « deux demonstrations de cette
généralisation », dont la seconde refait le chemin de l'article 357.

En termes actuels, le théorème de Gauss est le suivant. Pour $p$ premier
impair, soit $\Phi_p(x) = (x^p-1)/(x-1)$ le $p$-ième polynôme cyclotomique et
$p^{*} = (-1)^{(p-1)/2}p$. Alors il existe $Y, Z \in \mathbf{Z}[x]$ avec
\[
  4\,\Phi_p(x) = Y(x)^2 - p^{*}Z(x)^2 ,
\]
c'est-à-dire $Y^2 - pZ^2$ si $p \equiv 1 \pmod 4$ et $Y^2 + pZ^2$ si
$p \equiv 3 \pmod 4$.\footnote{Le signe que le brouillon laisse ouvert
(« $\pm$ ») est donc déterminé par $n$ modulo 4 ; l'énoncé est ici rendu plus
précis que le feuillet. C'est le contenu de l'article 357 des
\emph{Disquisitiones}, auquel elle renvoie et que la transcription cite : la
raison de fond en est que $\sqrt{p^{*}}$ appartient au corps cyclotomique
$\mathbf{Q}(\zeta_p)$ — la somme de Gauss quadratique —, ce qui est un langage
postérieur au feuillet.} Pour $p = 3$ on a $4\Phi_3 = (2x+1)^2 + 3$ ; pour
$p = 5$, $4\Phi_5 = (2x^2+x+2)^2 - 5x^2$.

Sa généralisation est vraie, et elle se démontre en une ligne à partir du
théorème lui-même. Puisque $x^{n^s} - 1 = \prod_{j=0}^{s}\Phi_{n^j}(x)$ avec
$\Phi_1(x) = x-1$, on a
\[
  \frac{x^{n^s}-1}{x-1} \;=\; \prod_{j=1}^{s}\Phi_{n^j}(x),
  \qquad \Phi_{n^j}(x) = \Phi_n\!\left(x^{n^{j-1}}\right),
\]
de sorte que chaque facteur, multiplié par 4, est déjà de la forme voulue :
$4\Phi_{n^j}(x) = Y\!\left(x^{n^{j-1}}\right)^2 - n^{*}Z\!\left(x^{n^{j-1}}\right)^2$.
Reste à multiplier ces $s$ formes entre elles, ce que fait l'identité de
composition
\[
  (Y^2 - n^{*}Z^2)(Y'^2 - n^{*}Z'^2)
  = (YY' + n^{*}ZZ')^2 - n^{*}(YZ' + ZY')^2 ,
\]
et à diviser par $4^{s-1}$. La division est licite : si $4$ divise $Y^2 -
n^{*}Z^2$ alors $Y \equiv Z \pmod 2$, et les deux composantes
$YY'+n^{*}ZZ'$ et $YZ'+ZY'$ sont alors paires, ce qui permet de sortir un
facteur 4 à chaque composition — il y en a $s-1$.\footnote{Cette démonstration
est de la présente lecture. Les deux démonstrations qu'elle annonce étaient
jointes à la lettre et ne sont pas dans le volume ; on ne sait donc pas si
c'est celle-ci, ni si les siennes étaient justes. Ce qui est établi ci-dessus
est que l'énoncé, lui, l'est.} On notera que le signe est le même pour toutes
les valeurs de $s$, puisque c'est toujours le même $n^{*}$.

\subsection*{Le passage de $Y, Z$ à $Y', Z'$}

Elle remarque ensuite que si l'on connaît $Y$ et $Z$ dans
$4(x^n-1)/(x-1) = Y^2 \pm nZ^2$ et qu'on veut $Y'$ et $Z'$ dans
$4(x^n+1)/(x+1) = Y'^2 \pm nZ'^2$, « il est clair qu'il suffiroit de changer
les signes de tous les termes de $Y$ et $Z$ qui contiennent des puissances
dont l'exposant est impair ».

C'est juste, et c'est la substitution $x \mapsto -x$. Pour $n$ impair,
\[
  \frac{x^n+1}{x+1} = \frac{(-x)^n - 1}{(-x) - 1},
\]
donc $\Phi_n(-x)$, et il suffit de poser $Y'(x) = Y(-x)$, $Z'(x) = Z(-x)$ —
ce qui est exactement changer le signe des termes de degré
impair.\footnote{L'hypothèse que $n$ est impair n'est pas écrite sur le
feuillet ; elle est indispensable, sans quoi $(x^n+1)/(x+1)$ n'est même pas un
polynôme. Le contexte la fournit — $n$ est chez elle un nombre premier impair
depuis le début du paragraphe —, mais l'énoncé tel qu'il est écrit ne la
porte pas.}

\subsection*{Fermat, pour l'exposant $n = p - 1$}

Le brouillon contient, une fois biffée d'un trait oblique et une fois récrite
presque mot pour mot quelques lignes plus bas, cette annonce :
l'impossibilité de $x^n + y^n = z^n$ en nombres entiers « n'a encore été
demontrée que pour $n = 3$ et $n = 4$ ; je crois etre parvenu a prouver cette
impossibilité pour $n = p-1$, $p$ etant un nombre premier de la forme
$8k+7$ ».\footnote{Le déplacement est significatif : biffée, la phrase ouvrait
le développement ; récrite, elle le ferme, comme « la derniere » des
« quelques autres considerations » ajoutées à l'article 357. C'est une
observation sur l'état du feuillet, non sur les mathématiques.}

Le premier membre de la phrase est exact pour l'époque : au début du siècle,
seuls les exposants 3 et 4 étaient acquis.

Le second demande à être précisé, car ce que la condition impose n'est pas
évident. Si $p \equiv 7 \pmod 8$, alors $n = p-1 \equiv 6 \pmod 8$ : $n$ est
pair, mais $4 \nmid n$, de sorte que le cas $n=4$ ne donne rien. Écrivons
$n = 2m$ avec $m = (p-1)/2$, qui est impair et $\equiv 3 \pmod 4$. L'équation
devient
\[
  (x^2)^m + (y^2)^m = (z^2)^m ,
\]
c'est-à-dire l'exposant impair $m$, restreint aux carrés. Et $p = 2m+1$ : le
nombre premier qu'elle choisit est précisément un nombre premier auxiliaire
pour l'exposant $m$, au sens de l'énoncé qu'on appelle aujourd'hui son
théorème.\footnote{Le nom « théorème de Sophie Germain » vient de
l'attribution que Legendre en fait en 1823 ; il n'est pas d'elle, et aucun
feuillet de ce volume n'en énonce le contenu. Le rapprochement fait ici est
celui de la présente lecture, à partir de la seule forme $p = 2m+1$.}

Dans cette situation, ce qu'on appelle le premier cas — $p \nmid xyz$ — est
immédiat et ne demande pas la condition $8k+7$ : si $p$ ne divise pas $x$,
alors $x^n = x^{p-1} \equiv 1 \pmod p$ par le petit théorème de Fermat, de
sorte que les trois termes valent $0$ ou $1$ modulo $p$, et que
$1 + 1 \equiv 1 \pmod p$ est impossible. Tout le contenu de son annonce est donc dans
le second cas, celui où $p$ divise l'un des trois nombres, et c'est là que la
condition sur $8k+7$ — laquelle équivaut à dire que 2 est résidu quadratique
et $-1$ non-résidu modulo $p$ — devait servir.\footnote{Ce n'est pas sans
rapport avec ce que Gauss loue dans sa lettre du 16 juin 1805 : « Votre
nouvelle demonstration pour les nombres premiers, dont 2 est residu ou non
residu m'a extremement plu ; elle est très fine, quoique elle semble être
isolée et ne pouvoir s'appliquer à d'autres nombres. » Le rapprochement est
de la présente lecture : le feuillet ne dit pas que les deux travaux sont
liés.}

La démonstration n'est pas dans le volume — elle était jointe à la lettre — et
rien ici ne permet de dire si elle valait. Deux remarques seulement, qui
tiennent au seul énoncé. Il est vide pour le plus petit $p$ de la forme
voulue : $p = 7$ donne $n = 6$, dont l'impossibilité suit du cas $n = 3$
qu'elle vient de nommer ; et il en va de même chaque fois que $3$ divise
$p-1$. Il n'a de contenu qu'à partir de $p = 23$, soit $n = 22$, puis
$p = 47$, soit $n = 46$. Et l'impossibilité pour tous ces exposants est
aujourd'hui acquise, comme conséquence du théorème démontré par Wiles en
1995 — deux siècles après le feuillet, et par des moyens qui n'ont aucun
rapport avec les siens.

\subsection*{La quantité de Lagrange (1775) : un calcul qui s'interrompt}

Le brouillon reproche ensuite à Lagrange, dans le mémoire de l'Académie de
Berlin pour 1775, de n'avoir pas su réduire la quantité de la page 352
\[
  s^{10} - 11\left(s^8 - 4s^6r^2 + 7s^4r^4 - 5s^2r^6 + r^8\right)r^2
\]
à la forme $t^2 - 11u^2$, et commence la réduction. Le calcul s'arrête sur une
parenthèse ouverte, au bout de trois lignes.

La présente lecture ne l'achève pas, et il faut dire pourquoi. La
transcription porte deux copies de la même quantité à deux lignes de
distance, et elles diffèrent : la seconde écrit $7s^4$ là où la première écrit
$7s^4r^4$. L'homogénéité tranche — la première lecture rend l'expression
homogène de degré 10, la seconde ne l'est pas —, mais elle ne dit pas laquelle
des deux copies est celle du feuillet et laquelle est une faute de lecture. La
transcription signale en outre que les exposants de la dernière ligne, et le
premier terme lu « $r^{10}$ », sont d'une lecture douteuse. Deux copies qui
divergent, une seule ligne de développement, et ses exposants incertains : il
n'y a pas là de calcul à vérifier.

Ce qu'on peut dire sans le feuillet : ramener une quantité à la forme
$t^2 - 11u^2$, c'est la représenter par la forme principale de discriminant
44, autrement dit l'écrire comme une norme de $\mathbf{Z}[\sqrt{11}]$. Et la
raison qu'elle donne de sa remarque est la bonne : « Cette remarque est une
nouvelle preuve de l'avantage de votre methode, qui, s'appliquant a toutes les
valeurs de $n$, donne pour chaque cas des valeurs de $Y$ et $Z$ independantes
de tatonnement. » Une méthode qui construit la représentation vaut mieux
qu'une méthode qui la cherche : c'est ce que la démonstration par composition
donnée plus haut fait aussi.

\subsection*{Formes ternaires, quaternaires, et le déterminant de l'adjointe}

Le second brouillon porte la plus nette de ses mathématiques. Elle rappelle
avoir réduit, suivant l'indication de Gauss, les formes ternaires de
déterminant nul aux formes binaires, et avoir cherché si la même chose vaut
pour les quaternaires ; puis elle énonce :

\begin{quote}
Je crois qu'en general $D$ etant la déterminante d'une forme composée d'un
nombre $n$ de variables, $D^{n-1}$ est la déterminante de l'adjointe de cette
forme.
\end{quote}

Elle observe que Gauss a trouvé $D^2$ pour l'adjointe ternaire, que son propre
calcul donne $D^3$ pour la quaternaire, et — c'est la phrase qui compte —
que « cette analogie n'est sans doute pas suffisante pour etablir la generalité
de la proposition ».

Les deux énoncés sont vrais, et tous deux sont aujourd'hui élémentaires.

Le premier : une forme quadratique dont la matrice de Gram est singulière
possède un radical non nul, et le quotient par ce radical est une forme non
dégénérée en un nombre strictement plus petit de variables. Une forme
quaternaire de déterminant nul se ramène donc à une forme en trois variables
au plus, comme elle le suppose — sur les rationnels sans réserve ; sur les
entiers, il faut de plus que le radical soit un facteur direct
primitif.\footnote{Cette réserve est de la présente lecture ; le feuillet
n'énonce que l'analogie avec le cas ternaire, et la question de l'équivalence
entière n'y est pas posée.}

Le second est l'identité
\[
  \det\left(\operatorname{adj} A\right) = \left(\det A\right)^{n-1}
\]
pour une matrice $n \times n$, qui se déduit immédiatement de
$A \cdot \operatorname{adj} A = (\det A)\,I$ en prenant les déterminants,
puis en traitant à part le cas $\det A = 0$ par continuité ou par densité des
matrices inversibles.\footnote{Gauss appelle « déterminante » d'une forme ce
que nous appelons le déterminant de sa matrice de Gram, à sa normalisation
près ; l'adjointe est notre comatrice transposée. La correspondance entre son
vocabulaire et le nôtre est faite ici une fois pour toutes.} Elle vérifie
$n = 3$ et $n = 4$ et refuse d'en conclure — ce qui est la bonne attitude,
puisque deux cas ne font pas une démonstration.

Elle donne pourtant un argument, et il faut voir exactement ce qu'il vaut : le
déterminant d'une forme à $n$ variables est un polynôme de degré $n$ en les
coefficients, ceux de l'adjointe sont de degré $n-1$, donc le déterminant de
l'adjointe est de degré $n(n-1)$ — « c'est a dire de l'ordre $n(n-1)$ » — et
$D^{n-1}$ l'est aussi. C'est un contrôle de degrés, c'est-à-dire une condition
nécessaire, exactement comme une vérification d'homogénéité : cela ne démontre
pas l'identité, mais cela l'aurait réfutée si les degrés avaient différé. Elle
ne prétend pas autre chose.

Le paragraphe suivant descend au cas binaire, et c'est là que le feuillet
résiste. La proposition donnerait, pour $n = 2$, que la déterminante de
l'adjointe est celle de la forme même — ce qui est correct, $D^{n-1} = D$ — et
que les coefficients de l'adjointe sont « de l'ordre $0$ ». Le sens demande
$n - 1 = 1$, non 0 ; et la transcription, qui donne le chiffre pour une lecture
douteuse, propose elle-même « $n-1$ » comme autre lecture possible. La présente
lecture retient ce que les mathématiques exigent sans effacer le
doute.\footnote{La transcription qualifie tout ce passage de « très remanié,
dont la lecture est incertaine », signale que « déterminante » est écrit
au-dessus d'un « l'adjointe » sans qu'on sache lequel des deux il remplace, et
que la première ligne est peut-être biffée. Aucune restitution du passage n'est
proposée ici au-delà du chiffre.}

\section*{Trois notes de lecture d'arithmétique}

\pagerange{91}{93}
Trois petits feuillets détachés, chacun en tête duquel elle écrit le numéro de
la page d'un livre qu'elle ne nomme pas.\footnote{Les feuillets 85 et 87
portent « p. 217 » et « p. 220 » ; le feuillet 86 renvoie à « Gauss pag. 272 ».
Que les trois renvoient au même ouvrage est l'hypothèse de la transcription
pour les deux premiers ; rien ne l'établit, et le troisième nomme Gauss.}

\subsection*{$p^4 + 4$}

« Aucun nombre de la forme $p^4+4$ excepté 5 n'est un nombre premier, car
$p^4+4 = (p^2-2)^2 + 4p^2$ et par conséquent ces nombres sont, de plusieurs
manières, la somme de deux quarrés. » Elle ajoute que pour $p = 1$,
$(p^2-2)^2 = 1$ et « les deux membres sont identiques ».

L'énoncé est vrai. Son argument repose sur deux représentations du même
nombre,
\[
  p^4 + 4 = (p^2)^2 + 2^2 = (p^2-2)^2 + (2p)^2 ,
\]
et sur le théorème d'Euler suivant lequel un entier admettant deux
représentations essentiellement différentes comme somme de deux carrés est
composé. Il faut, pour l'appliquer, que les deux paires soient distinctes, et
elles ne le sont pas toujours : elles coïncident pour $p = 1$, où l'on trouve
$5 = 1 + 4$ — c'est l'exception qu'elle relève elle-même — mais aussi pour
$p = 2$, où les deux paires sont $\{4,2\}$ et $\{2,4\}$. Ce second cas ne
menace pas la conclusion, puisque $p^4+4$ est alors pair, mais la restriction
manque au feuillet.\footnote{C'est l'un des quatre défauts qu'une lecture
modernisée doit chercher : une hypothèse tacite. Le feuillet écarte $p=1$ et
ne dit rien de $p=2$.}

La factorisation explicite que son argument n'exhibe pas est celle qu'on
appelle aujourd'hui l'identité de Sophie Germain :
\[
  a^4 + 4b^4 = \left(a^2 + 2b^2 - 2ab\right)\left(a^2 + 2b^2 + 2ab\right),
\]
soit, pour $b = 1$, $p^4 + 4 = (p^2-2p+2)(p^2+2p+2)$. Elle règle les deux cas
d'un coup : le nombre n'est premier que si le premier facteur vaut 1,
c'est-à-dire si $p = 1$.\footnote{Le feuillet ne porte pas cette identité : il
porte la démonstration par les deux sommes de deux carrés, qui est autre. Le
nom, comme tous les noms modernes, lui est postérieur ; la présente lecture ne
se prononce pas sur la question de savoir d'où il vient, laquelle relève des
sources et non des feuillets.}

\subsection*{$p^4 + q^4$ et les nombres de Fermat}

Elle écrit ensuite
\[
  p^4 + q^4 = \left(p^2 - q^2\right)^2 + 2(pq)^2 = \left(p^2 + q^2\right)^2 - 2(pq)^2 ,
\]
deux identités exactes, et en tire que « lorsque $p^4+q^4$ est un nombre
premier on peut être assuré que ce nombre n'est que de ces seules manières la
somme d'un quarré $+$ ou $-$ un double quarré ».

C'est encore juste, et la raison moderne en est que les anneaux
$\mathbf{Z}[\sqrt{-2}]$ et $\mathbf{Z}[\sqrt{2}]$ sont principaux : un nombre
premier y a, à unités et conjugaison près, une seule décomposition, de sorte
que sa représentation par $x^2 + 2y^2$, comme sa représentation par
$x^2 - 2y^2$, est unique. Les deux représentations qu'elle écrit sont donc
\emph{les} représentations.

Elle finit par : « Les nombres $2^{2^i}+1$ ne sont qu'un cas particulier de la
forme $p^4+q^4$. » C'est vrai à partir de $i = 2$, avec
$p = 2^{2^{i-2}}$ et $q = 1$ : ainsi $17 = 2^4 + 1^4$,
$257 = 4^4 + 1^4$, $65537 = 16^4 + 1^4$. Ce ne l'est pas pour $i = 0$ et
$i = 1$, qui donnent 3 et 5.\footnote{L'hypothèse $i \geq 2$ est implicite sur
le feuillet, qui écrit la remarque sans restriction. C'est le genre de
restriction qu'une lecture modernisée doit rétablir plutôt que reconduire.}
La remarque est utile : elle place les nombres de Fermat dans une famille où
les deux identités ci-dessus s'appliquent, et donc où l'unicité de la
représentation par $x^2 \pm 2y^2$ fournit un test.

\subsection*{L'objection à Gauss sur la composition des formes binaires}

Le feuillet 86 (vue 92) est le plus intéressant du lot, et le plus discutable.
Il tient en quatre phrases :

\begin{quote}
Gauss pag. 272 dit que le produit des deux formes $(16, 3, 19)$, $(8, 1, 37)$
est $(8, -1, 37)$. Ce produit est aussi $(13, -2, 23)$, on peut s'en assurer.
En général il y a toujours deux formes qui résultent du produit de deux formes
données, mais dans les cas particuliers les deux formes se réduisent
quelquefois à une. Il semble que l'auteur n'a pas fait cette remarque, et alors
il y auroit plusieurs choses à rectifier dans la théorie qu'il donne du produit
de plusieurs formes quadratiques binaires.
\end{quote}

Prenons les quatre formes au sérieux.\footnote{Gauss note $(a,b,c)$ la forme
$ax^2 + 2bxy + cy^2$ et appelle « déterminante » le nombre $b^2 - ac$ ; le
discriminant d'aujourd'hui en est le quadruple. La transcription conserve sa
notation ; elle est conservée ici aussi, faute de quoi la comparaison avec le
feuillet deviendrait illisible. Les douze nombres sont donnés tels que la
transcription les lit, sans vérification sur l'image.} Les quatre ont le même
déterminant : $3^2 - 16\cdot 19 = 1 - 8\cdot 37 = (-1)^2 - 8\cdot 37 =
(-2)^2 - 13\cdot 23 = -295$. Elles appartiennent donc bien au même ordre, où
la composition a un sens. Les formes réduites de ce déterminant dans l'ordre
où Gauss compose sont au nombre de huit :
\[
  (1,0,295),\quad (5,0,59),\quad (8,\pm 1,37),\quad (13,\pm 2,23),\quad
  (16,\pm 3,19),
\]
et les deux formes qu'elle nomme, $(8,-1,37)$ et $(13,-2,23)$, sont deux
d'entre elles — deux formes réduites distinctes, donc deux classes distinctes,
car deux formes réduites différentes ne sont jamais proprement
équivalentes.\footnote{Cette énumération et les vérifications qui suivent sont
de la présente lecture ; le feuillet ne porte aucun calcul, seulement
l'assertion « on peut s'en assurer ».}

Le calcul donne raison à Gauss. La composée de $(16,3,19)$ et de $(8,1,37)$
est bien $(8,-1,37)$. Et la seconde forme qu'elle propose n'est pas une
seconde valeur du même produit : c'est la composée de $(16,-3,19)$ — la forme
\emph{opposée} de la première — avec $(8,1,37)$ :
\[
  (16,3,19) \circ (8,1,37) = (8,-1,37),
  \qquad
  (16,-3,19) \circ (8,1,37) = (13,-2,23).
\]
Autrement dit, les deux formes qu'elle obtient diffèrent exactement du
changement de signe du coefficient médian de l'un des deux facteurs,
c'est-à-dire du passage d'une forme à son opposée ; et les deux résultats
diffèrent alors de la classe ambiguë $(5,0,59)$, qui est d'ordre 2 dans le groupe des
classes.\footnote{On vérifie de même que
$(13,-2,23) = (8,-1,37) \circ (5,0,59)$.}

Ce qui reste de son objection est donc ceci, et c'est réel : si l'on compose
sans avoir fixé l'orientation des deux formes de départ — si l'on ne distingue
pas $(a,b,c)$ de $(a,-b,c)$, qui représentent les mêmes nombres et sont
équivalentes par une substitution de déterminant $-1$ —, alors le produit
n'est pas déterminé, et l'on obtient deux formes. Ce qui ne reste pas, c'est
la conclusion : Gauss a fait la remarque, et c'est précisément pour cela qu'il
impose à sa transformation les conditions de signe de l'article 235. Ce sont
elles qui font que la composition passe au quotient et définit une opération
sur les classes pour l'équivalence \emph{propre} — le groupe des classes, qui
compte ici huit éléments. Rien n'est donc à rectifier dans
sa théorie ; ce qui manque au feuillet, c'est la distinction entre équivalence
propre et impropre, qui est le pivot de la construction.

Il faut ajouter la réserve d'usage : les douze nombres du feuillet n'ont pas été
vérifiés sur l'image, et la démonstration ci-dessus est suspendue à leur
exactitude. Que les quatre formes aient toutes le même déterminant, ce qui
n'était nullement garanti d'avance, est un indice sérieux qu'elles sont bien
lues.

\section*{La note de Libri}

\pagerange{65}{65}
Un feuillet d'une autre main, signé « G. Libri » et daté 1824 (f. 60), porte
trois congruences sur la décomposition d'un nombre premier $p = 4k+1$ en
$a^2+b^2$, $a$ impair et $b$ pair. La dernière, qui est celle qui porte la
conclusion, est un théorème classique :
\[
  \frac{2}{1}\cdot\frac{6}{2}\cdot\frac{10}{3}\cdots\frac{4k-2}{k}
  \;\equiv\; \pm\,2a \pmod{p},
\]
avec le signe $+$ si $a \equiv 1 \pmod 4$ et $-$ si $a \equiv -1 \pmod 4$. Le
produit de gauche est le coefficient binomial central : en effet
$\prod_{j=1}^{k}(4j-2)/j = 2^k(2k-1)!!/k! = \binom{2k}{k}$, de sorte que
l'énoncé est
\[
  \binom{2k}{k} \equiv 2a \pmod p, \qquad a \equiv 1 \pmod 4 ,
\]
qui est le résultat de Gauss sur le coefficient binomial central. On le
vérifie sur les premiers cas : $p=13$, $k=3$, $\binom63 = 20 \equiv 7$ et
$a = -3$, $2a = -6 \equiv 7$ ; $p=17$, $\binom84 = 70 \equiv 2 = 2a$ avec
$a = 1$ ; $p=29$, $\binom{14}{7} = 3432 \equiv 10 = 2a$ avec $a=5$.

La conclusion qu'il en tire — « On peut décomposer de cette manière
directement le nombre $p$ en deux carrés » — est juste, et mérite d'être dite
en clair : comme $a^2 < p$, on a $|2a| < 2\sqrt{p}$, de sorte que le résidu de
$\binom{2k}{k}$ modulo $p$ détermine $a$ sans ambiguïté dès que
$4\sqrt{p} < p$, c'est-à-dire dès que $p > 16$ ; et $b$ s'obtient alors par
$b = \sqrt{p - a^2}$. C'est bien un algorithme de décomposition, et non une
simple identité.

La deuxième congruence du feuillet,
\[
  \left[(k+1)(k+2)\cdots 2k\right]^2 \equiv \pm\,2b \pmod p,
\]
est vraie elle aussi, et se déduit de la précédente. Le produit vaut
$\left(\frac{p-1}{2}\right)!\,/\,k!$ ; comme
$\left(\frac{p-1}{2}\right)!^{\,2} \equiv -1 \pmod p$ pour $p \equiv 1
\pmod 4$, son carré vaut $-1/(k!)^2$, c'est-à-dire
$-2a\,/\,\left(\frac{p-1}{2}\right)!$ d'après la congruence précédente ; et
comme $\left(\frac{p-1}{2}\right)!$ et $a/b$ sont l'une et l'autre des
racines carrées de $-1$ modulo $p$, elles ne diffèrent que du signe, ce qui
laisse $\mp 2b$. On le vérifie : pour $p = 5$ le carré vaut $4 \equiv 2b$ avec
$b = 2$ ; pour $p = 13$, $3^2 = 9 \equiv -4$ avec $b = 2$ ; pour $p = 17$,
$14^2 \equiv 9 \equiv -8$ avec $b = 4$ ; pour $p = 29$,
$27^2 \equiv 4 = 2b$ avec $b = 2$. La première, en revanche — la même sans le
carré, avec le même second membre $\pm 2b$ — est fausse : pour $p = 5$ elle
demanderait $2 \equiv \pm 4$, et pour $p = 13$, $120 \equiv 3 \equiv \pm 4$,
ce qui n'a lieu ni dans un cas ni dans l'autre.\footnote{La transcription
signale précisément cette anomalie — « les deux premières congruences portent
le même second membre, bien que la seconde soit le carré de la première : la
feuille est ainsi, et rien n'y est biffé » — sans la résoudre. La présente
lecture ne la résout pas davantage : elle constate que la seconde est vraie,
que la première ne l'est pas, et que la faute peut être du feuillet comme de
la lecture du feuillet.}

\section*{Le levier, l'infini et l'hyperbole}

\pagerange{40}{42}
Trois vues portent un texte d'une main qui n'est ni la sienne ni celle
d'aucun correspondant identifié, intitulé « Réponse aux questions de Mlle
\ldots{} » — le nom est réduit à un trait ondulé. C'est un cours, écrit pour
quelqu'un qui apprend, et il vaut d'être lu comme tel : il donne la mesure de
ce qu'on lui enseignait.

Le problème est celui du levier avec trois puissances $P$, $S$, $R$ appliquées
en trois points $A$, $C$, $B$ d'une droite et en équilibre ; l'auteur établit
$S \cdot CB = P \cdot AB$ en prenant $B$ pour point d'appui, puis fait
s'éloigner $P$ : à distance cent fois plus grande, une force cent fois moindre
suffit, et « à une distance infiniment grande une force infiniment petite ou
nulle suffira ». Le paradoxe apparent qui résulte, dit l'auteur, d'une
proposition de Monge — qu'une droite rigide infinie, heurtée en un point, ne
peut prendre de mouvement — est alors levé par un argument
qui est encore le nôtre : « pour que la ligne infinie $BA$ prît quelque
mouvement, il faudroit que le point $A$ parcourût dans un temps fini un chemin
infiniment grand ; ce qui est absurde ».

Suit la définition que l'auteur donne de l'infini, et qui est la seule phrase
du volume qui soit de la logique : « dans le langage rigoureux de la
geometrie, l'infini n'est qu'une grandeur plus grande que toute grandeur
donnée ». C'est la définition d'aujourd'hui, celle qui refuse à l'infini le
statut d'objet et n'en fait qu'une manière de parler d'une famille de
grandeurs. Tout le reste du morceau en découle. L'hyperbole entre ses
asymptotes fournit le paradoxe géométrique correspondant : le rectangle
d'ordonnée $Ee$ et d'abscisse $Ce$ est constant, égal au carré $CB$, quand
$Ce \to \infty$ et $Ee \to 0$ ; et comme on peut prendre une autre hyperbole
entre les mêmes asymptotes, « $0 \times \infty$ » vaudrait aussi bien n'importe
quelle autre grandeur finie. De là, en substituant $\infty = 1/0$, la
proposition que $0/0$ est « egal a tout ce qu'on veut ». La conclusion de
l'auteur est exacte : « ces paradoxes n'ont lieu que parceque c'est un
paradoxe lui meme que de supposer qu'on puisse atteindre l'infini ».

En termes actuels, ce que le morceau décrit est une forme indéterminée : la
limite d'un produit $f(x)g(x)$ où $f \to 0$ et $g \to \infty$ n'est pas
déterminée par ces deux données seules, et l'hyperbole en donne la famille
d'exemples la plus simple, chaque hyperbole $xy = c$ réalisant une valeur
$c$.\footnote{La reformulation par les limites est de la présente lecture ; le
feuillet n'a ni le mot ni la notion, et son argument est celui, plus ancien et
parfaitement valable, du refus d'un infini actuel.}

\section*{Le centre de gravité du triangle}

\pagerange{98}{99}
Deux feuillets de sa main, l'un une première ébauche (f. 84), l'autre une
rédaction plus complète (ff. 92r–92v), traitent le même problème sous le
renvoi « fig 39 » : trouver le centre de gravité d'un triangle. La méthode est
une méthode d'échelle, et elle est élégante.

Le triangle $ABC$ est coupé par les milieux en un parallélogramme $AFDE$ et
deux triangles $BFD$, $CED$ semblables à $ABC$ dans le rapport $\frac12$. Posons
$a = \frac12 AD$, mesurons les distances à la base $BC$ le long de la médiane
$AD$, et appelons $2x$ la distance cherchée du centre de gravité de $ABC$ à
$BC$. Par similitude, les centres de gravité des deux petits triangles sont à
la distance $x$ de leurs bases, donc de $BC$ ; le centre du parallélogramme
est le milieu $C'$ de $AD$, à la distance $a$. Les deux parts ont même aire —
le parallélogramme fait la moitié du triangle, les deux petits triangles
l'autre moitié — de sorte que le centre de gravité de l'ensemble est le milieu
des deux :
\[
  2x = \frac{a + x}{2}, \qquad 4x = a + x, \qquad x = \frac{a}{3},
\]
et le centre de gravité est à la distance $2x = \frac{2a}{3} = \frac{AD}{3}$
de la base : il partage la médiane dans le rapport $2 : 1$. C'est le résultat
juste, et c'est celui qu'elle écrit, « c. q. f. d. », dans la marge de la vue
99.\footnote{Le feuillet écrit « $C'X = \frac12(a-x)$ » là où la construction
donne $C'X = a - x$ ; c'est l'expression suivante, « $x + \frac12(a-x)$ », qui
est la bonne, et qui vaut bien $(a+x)/2$. La ligne se perd hors du cadre à
droite, et la transcription signale que la colonne marginale porte encore deux
ou trois lignes à demi biffées et illisibles. La chaîne complète du calcul est
donc restituée ici à partir de ses deux extrémités — l'énoncé du feuillet et
sa conclusion marginale $4x = a + x$ —, et le pas intermédiaire n'est pas lu
mais reconstruit.}

La première ébauche (f. 84) contient une inadvertance que la transcription
relève : elle écrit deux fois « $EE'$ » là où le sens demande « $FF'$ » pour la
seconde des deux médianes.

\section*{Ce qui n'est pas des mathématiques}

\pagerange{7}{89}
Plus de la moitié du volume. Ces pièces ne sont pas résumées ici pour ce
qu'elles disent des mathématiques — elles n'en disent rien — mais parce qu'un
lecteur qui ouvre Français 9118 les rencontrera d'abord, et qu'aucune n'est un
défaut de la transcription.

\subsection*{Les lettres d'académie}

\pagerange{12}{22}
Delambre, secrétaire perpétuel, lui envoie deux billets pour la séance
publique du 6 janvier 1816 et souhaite qu'elle s'y rende elle-même ; il accuse
réception, le 23 juillet 1821, des \emph{Recherches sur la théorie des surfaces
élastiques} déposées à la bibliothèque de l'Institut. Fourier lui écrit le 30
mai 1823 qu'elle sera désormais admise « dans l'une des places reservées au
centre de la salle », distinction que l'Académie accorde en considération de
ses travaux et du grand prix qu'elle a obtenu. Cauchy, le 24 juillet 1821,
accuse réception d'un ouvrage qu'elle lui a adressé et lui envoie en retour son
\emph{Analyse algébrique} ; le 23 juillet 1826, d'un mémoire. Navier, le 2 août 1821, écrit la
phrase du volume qu'on citera : « un écrit aussi remarquable, que bien peu
d'hommes peuvent lire, et qu'une seule femme pouvait faire ». Les billets
autographes de Fourier sont d'un autre ton — des places de loge aux Italiens,
un rhume qu'il appelle « une courbature » par plaisanterie de géomètre, et
l'élection de novembre où il compte les voix de Legendre, de Jussieu et de
Montmorency.

\subsection*{Les récépissés de l'Institut}

\pagerange{85}{89}
Trois feuilles de formulaire, signées Cardot, et les seules pièces datées avec
certitude de tout le dossier du concours. Le 21 septembre 1811 et le 21
septembre 1813, un mémoire est reçu pour le prix sur les vibrations des
surfaces élastiques, avec une épigraphe de Newton la première fois, de Bacon
la seconde ; le 8 novembre 1821, un mémoire est reçu pour le prix de
mathématiques de 1822, sous une épigraphe de Gauss. Chacun des trois porte
« que j'ai numéroté I » — ce qui s'accorde avec la lettre de Legendre du 22
octobre 1811 : « Votre mémoire n'est pas perdu : il est le seul qu'on ait reçu
sur la question des vibrations des surfaces. »

\subsection*{Gauss vu de Brunswick, novembre 1806}

\pagerange{68}{71}
Deux lettres qui ne contiennent pas une ligne de mathématiques, et qui portent
pourtant l'épisode le plus concret du volume. Elle avait chargé le général
Pernety de s'informer de Gauss, l'armée française occupant Brunswick ; Pernety
chargea un officier, Chantel, lequel se rendit chez Gauss, le trouva bien portant, « un
peut Confus » de ne connaître ni le général ni la demoiselle, le recommanda au
gouverneur de la place et l'emmena dîner. Pernety lui transmet la lettre
depuis les lignes du siège de Cosel, en décembre, avec une phrase où l'on
entend son métier : « faisant par Réflexion tout le mal que je peux à qui
jamais ne m'en fit aucun et que je ne connois pas, mais c'est le métier ». Il
appelle Gauss « cet émule d'Archimède, mieux traité que lui ».

\subsection*{Villoison, les vers, et la plaquette}

\pagerange{76}{85}
Trois lettres de d'Ansse de Villoison — des 12 et 14 juillet 1802, la
troisième sans date —, et la plaquette imprimée à laquelle elles se
rapportent. L'helléniste avait composé, pour
l'anniversaire de Lalande, des vers latins où une note en bas de page nommait
« Mademoiselle Sophie Germain » et disait qu'elle excelle dans les
mathématiques ; elle exigea qu'il se rétracte. Il répond qu'il a brûlé sa pièce
de vers grecs, qu'il voudrait pouvoir anéantir de même les latins, et donne
« ma parole d'honneur que jamais je ne me permettrai de parler de vous \ldots{}
dans aucun écrit ». L'exemplaire de la plaquette conservé dans le volume porte
une correction manuscrite à l'encre sur un vers de la page (7). Elle avait donc
gardé, dans ses papiers, la pièce même qu'elle avait fait retirer.

\subsection*{Trois feuillets de physique}

\pagerange{94}{96}
Trois feuillets d'une écriture plus ronde et plus posée que celle de ses
billets de calcul, sur des sujets de physique : un raisonnement de « Monsieur
Abeille » suivant lequel le froid serait un être réel, puisque l'eau augmente
de volume en gelant et que le néant ne produit aucun effet positif ; une
réponse à Abeille sur le feu et la lumière, où l'auteur soutient que
l'impossibilité de décomposer le feu tient à l'impossibilité de l'obtenir seul
et non à la simplicité de son essence ; et des « Idées sur les Sens », qui
ramènent les cinq sens au toucher, l'œil étant sensible « à l'attouchement du
subtile élément de la lumière ». Ce sont des raisonnements de physique
qualitative du XVIII\textsuperscript{e} siècle, sans mathématiques ; ils sont
signalés ici et non modernisés, parce qu'il n'y a rien à y
moderniser.\footnote{La transcription range ces trois feuillets parmi les
feuillets détachés de sa main, puis note que l'écriture en est plus ronde et
plus posée que celle des billets de calcul. Le « mon opinion » de la seconde pièce
fait de l'auteur du raisonnement celui du feuillet, mais le volume ne le
nomme pas et la présente lecture ne l'attribue pas.}

\subsection*{Le reste}

\pagerange{74}{75}
Un billet du libraire Bernard, quai des Augustins, demandant à présenter « le
citoien Cousin » à sa mère (f. 3r–3v, vues 7 et 8) ; et une invitation à dîner
d'un correspondant anonyme du Cloître Notre-Dame, dont la signature est un
paraphe illisible, promettant du beurre
frais et des betteraves et une parente « qui ne connoît de géométrie que les
figures les plus naturelles ».\footnote{Le billet est daté « ce 17 pluviôse »
sans millésime. La transcription observe que « duodi » y désigne le dimanche
suivant, que le 22 pluviôse ne tombe un dimanche, dans ces années-là, qu'en
l'an VII, et en conclut le 5 février 1799 ; elle donne cela pour son
inférence et non pour une donnée du feuillet.}

\section*{Ce que les transcriptions laissent en suspens}

\pagerange{47}{64}
Cinq questions, que la présente lecture n'a pas pu trancher et n'a pas
tranchées.

\begin{itemize}
  \item \emph{La date de la lettre de Legendre du 19 janvier} (ff. 52r–53r).
    Le millésime se lit « 1814 », mais le feuillet d'adresse au verso porte un
    timbre de janvier 1811, la note à laquelle la lettre renvoie a été postée
    en janvier 1811, et 1811 est l'année du premier concours. La transcription
    donne les deux lectures et laisse la question ouverte ; elle le reste.
  \item \emph{Laquelle des notes de Legendre est « ma 1ère note »}. La lettre
    du 19 janvier dit revenir sur « la Multiplication par $\sin\frac12\omega$,
    contre laquelle je m'étois élevé dans ma 1ère note ». La transcription du
    lot 3 identifie cette première note à celle du § 47 (ff. 42r–43r) — mais
    cette note-là ne parle nulle part de $\sin\frac12\omega$, tandis que la
    note « Pag. 155 » (ff. 58–59), transcrite dans un autre lot, ne parle que
    de cela. Les deux lots ont été transcrits par deux passes différentes, et
    aucune n'a vu le feuillet de l'autre. La question de l'ordre des trois
    notes reste donc entière, et avec elle celle de savoir si Legendre a
    d'abord objecté puis concédé, ou l'inverse.
  \item \emph{Une contradiction entre deux de ces notes}. Sous
    $\sin\frac12\omega = 0$, la note « Pag. 155 » affirme que les coefficients
    $\gamma$, $\gamma'$, $\delta'$ deviennent infinis ; la lettre du 19 janvier
    affirme que $\delta'$ est nul et que seuls $\gamma$ et $\gamma'$
    survivent, indéterminés. Sans le texte d'Euler auquel les deux renvoient —
    que le volume ne contient pas — il n'y a pas moyen de dire laquelle a
    raison.
  \item \emph{Le destinataire du brouillon « ce 18 juillet »}, que le feuillet
    ne nomme pas.
  \item \emph{Ce qui manque}. Trois pièces essentielles sont annoncées dans le
    volume et n'y sont pas : les deux démonstrations de la généralisation du
    théorème de l'article 357, les « quelques autres considerations » où
    figurait l'argument sur l'équation de Fermat, et le mémoire de 1811
    lui-même, dont on n'a ici que le jugement. Le volume est le dossier d'un
    travail, non le travail.
\end{itemize}

Deux points de moindre portée, enfin, où le feuillet résiste et où la présente
lecture a écrit ce que les mathématiques exigent en le disant : le coefficient
de la première équation de la note de Lagrange, illisible, et le chiffre
« de l'ordre 0 » du second brouillon à Gauss, qui doit se lire 1.

\end{document}
